\paragrafos{Πολλαπλασιασμός κλασμάτων: το «του» είναι «επί»}{Α.2.5, «Πολλαπλασιασμός κλασμάτων»}

\begin{anagke}
Η Αγγελική ετοιμάζει ένα πίνακα για τη σχολική έκθεση. Παίρνει ένα ορθογώνιο χαρτόνι και ζωγραφίζει θάλασσα στα τρία τέταρτα του. Από τη θάλασσα, τα δύο τρίτα τα βάφει σκούρο μπλε, για το βάθος. «Τι μέρος του χαρτονιού είναι σκούρο μπλε;» ρωτάει η Ζωή.

Ο Ηλίας, που δοκιμάζει πάντα πρώτα κάτι με τα χέρια του, κόβει ένα χαρτί και το χωρίζει σε 4 ίσες στήλες. Σκιάζει τις 3. Μετά χωρίζει όλο το σχήμα σε 3 ίσες σειρές και σκουραίνει τις 2 από τις 3 σειρές πάνω στη σκιασμένη περιοχή. Μετράει: από 12 ίσα τετράγωνα, τα 6 είναι σκούρα. «Τα μισά!» λέει.

Η Μυρτώ σκέφτεται δυνατά: «Τα δύο τρίτα των τριών τετάρτων είναι το μισό του χαρτονιού. Αλλά όταν πολλαπλασιάζεις, δεν βγαίνει πάντα κάτι μεγαλύτερο; Εδώ βγήκε μικρότερο και από τα δύο!» Και η Πολυξένη προσθέτει: «Και αν βρούμε έναν κανόνα, ώστε να μην χρειάζεται να ζωγραφίζουμε κάθε φορά;»
\end{anagke}

\begin{lexi}{αντίστροφος}{αντίστροφος}
Η λέξη \lexiword{αντίστροφος} έρχεται από το «αντί» και το «στρέφω», δηλαδή «γυρίζω ανάποδα». Ο αντίστροφος ενός κλάσματος είναι το ίδιο κλάσμα γυρισμένο ανάποδα: ο αριθμητής γίνεται παρονομαστής και ο παρονομαστής αριθμητής. Δύο αριθμοί είναι αντίστροφοι όταν το γινόμενό τους είναι 1.

Και μια μικρή λέξη που κρύβει τα πάντα: το «\lexiword{του}» σε φράσεις όπως «τα $\frac{2}{3}$ του $\frac{3}{4}$» σημαίνει \emph{επί}. «Τα δύο τρίτα του» είναι «δύο τρίτα επί».
\end{lexi}

\begin{metafrasi}
Τα $\frac{2}{3}$ των $\frac{3}{4}$. & $\dfrac{2}{3}\cdot\dfrac{3}{4}=\dfrac{6}{12}=\dfrac{1}{2}$\\
Τα $\frac{3}{5}$ του 20. & $\dfrac{3}{5}\cdot 20=12$\\
Ο αντίστροφος του $\frac{3}{7}$. & $\dfrac{7}{3}$, γιατί $\dfrac{3}{7}\cdot\dfrac{7}{3}=1$\\
Ο αντίστροφος του 5. & $\dfrac{1}{5}$, γιατί $5\cdot\dfrac{1}{5}=1$\\
Δύο ολόκληρα και μισό, τέσσερις φορές. & $2\dfrac{1}{2}\cdot 4=\dfrac{5}{2}\cdot 4=10$\\
\end{metafrasi}

\tmima{Η θεωρία}

\begin{kanonas}{γινόμενο κλασμάτων}
Το γινόμενο δύο κλασμάτων είναι κλάσμα με αριθμητή το γινόμενο των αριθμητών και παρονομαστή το γινόμενο των παρονομαστών:
$$\dfrac{\alpha}{\beta}\cdot\dfrac{\gamma}{\delta}=\dfrac{\alpha\cdot\gamma}{\beta\cdot\delta}.$$
Το γινόμενο φυσικού επί κλάσμα είναι κλάσμα με αριθμητή το γινόμενο του αριθμητή επί τον φυσικό και τον ίδιο παρονομαστή: $\lambda\cdot\dfrac{\alpha}{\beta}=\dfrac{\lambda\cdot\alpha}{\beta}$.
\end{kanonas}

Γιατί δουλεύει; Πάρε το χαρτόνι της Αγγελικής. Χωρίζεται σε $4\cdot 3=12$ ίσα τετράγωνα, και τα σκούρα είναι $3\cdot 2=6$. Δηλαδή $\frac{3}{4}\cdot\frac{2}{3}=\frac{3\cdot 2}{4\cdot 3}=\frac{6}{12}$. Κάθε φορά οι παρονομαστές δίνουν πόσα τετράγωνα έχει το όλο και οι αριθμητές πόσα χρωματίστηκαν.

\begin{center}
\begin{tikzpicture}[x=0.9cm,y=0.9cm]
\fill[gblue!25] (0,0) rectangle (3,3);
\fill[gblue!70] (0,1) rectangle (3,3);
\foreach \i in {0,...,4}{\draw (\i,0)--(\i,3);}
\foreach \j in {0,...,3}{\draw (0,\j)--(4,\j);}
\draw[very thick] (0,0) rectangle (4,3);
\node[below] at (2,-0.1) {\small 4 στήλες, σκιασμένες 3};
\node[right] at (4.2,1.5) {\small 3 σειρές, σκούρες 2};
\end{tikzpicture}
\end{center}

\tmimasmall{Αντίστροφοι}

Δύο κλάσματα λέγονται \textbf{αντίστροφα} όταν το γινόμενό τους είναι 1: $\dfrac{\alpha}{\beta}\cdot\dfrac{\beta}{\alpha}=\dfrac{\alpha\beta}{\beta\alpha}=1$. Ο αντίστροφος του $\frac{3}{7}$ είναι ο $\frac{7}{3}$. Ο αντίστροφος του φυσικού $5=\frac{5}{1}$ είναι ο $\frac{1}{5}$. Ο αντίστροφος του 1 είναι ο ίδιος ο 1. Το 0 \textbf{δεν έχει αντίστροφο}, γιατί δεν υπάρχει αριθμός που επί 0 να κάνει 1.

Όλες οι ιδιότητες του πολλαπλασιασμού των φυσικών ισχύουν και στα κλάσματα: αντιμεταθετική, προσεταιριστική, επιμεριστική και ο αριθμός 1 που δεν αλλάζει το γινόμενο.

\tmimasmall{Μεικτοί και απλοποίηση}

Για να πολλαπλασιάσω μεικτούς αριθμούς, τους γράφω πρώτα ως κλάσματα: $2\frac{1}{3}\cdot 1\frac{1}{2}=\frac{7}{3}\cdot\frac{3}{2}$. Πριν πολλαπλασιάσω, μπορώ να απλοποιήσω «χιαστί»: το 3 του παρονομαστή και το 3 του αριθμητή απλοποιούνται, και μένει $\frac{7}{2}=3\frac{1}{2}$. Έτσι δουλεύω με μικρότερους αριθμούς.

\begin{prosoxi}
\begin{itemize}[leftmargin=1.2em,itemsep=0.3ex,topsep=2pt]
\item Ο πολλαπλασιασμός με κλάσμα μικρότερο από το 1 \textbf{μικραίνει}. Τα $\frac{2}{3}$ του $\frac{3}{4}$ είναι μικρότερα και από τα δύο: $\frac{1}{2}<\frac{2}{3}$ και $\frac{1}{2}<\frac{3}{4}$. Ο πολλαπλασιασμός με κλάσμα μεγαλύτερο από το 1 μεγαλώνει.
\item Δεν πολλαπλασιάζουμε χωριστά τα ολόκληρα και τα κλάσματα ενός μεικτού: $2\frac{1}{2}\cdot 3\frac{1}{2}$ δεν είναι $(2\cdot 3)+\left(\frac{1}{2}\cdot\frac{1}{2}\right)=6\frac{1}{4}$. Σωστά: $\frac{5}{2}\cdot\frac{7}{2}=\frac{35}{4}=8\frac{3}{4}$.
\item Δεν πολλαπλασιάζουμε «και παρονομαστές» όταν προσθέτουμε: στο άθροισμα κρατάμε τον ίδιο παρονομαστή, στο γινόμενο τους πολλαπλασιάζουμε.
\end{itemize}
\end{prosoxi}

\tmima{Λυμένα παραδείγματα}

\paradeigma{Το χαρτόνι της Αγγελικής}
\tixero Τα $\frac{3}{4}$ του χαρτονιού είναι θάλασσα, και τα $\frac{2}{3}$ της θάλασσας είναι σκούρα.\par
\tithelo Τι μέρος του χαρτονιού είναι σκούρο.\par
\possindeo Το «τα $\frac{2}{3}$ των $\frac{3}{4}$» σημαίνει $\frac{2}{3}\cdot\frac{3}{4}=\frac{2\cdot 3}{3\cdot 4}=\frac{6}{12}$. Το $\frac{6}{12}$ απλοποιείται με το 6: $\frac{1}{2}$. Μπορώ να απλοποιήσω και πριν: $\frac{2\cdot\cancel{3}}{\cancel{3}\cdot 4}=\frac{2}{4}=\frac{1}{2}$.\par
\apantisi Το $\dfrac{1}{2}$ του χαρτονιού.\par
\elegxos Το σχήμα της θεωρίας δείχνει 6 σκούρα τετράγωνα από 12. Σωστά.

\paradeigma{Οι μαθητές της Α’}
\tixero Η Α’ τάξη έχει 96 μαθητές και τα $\frac{5}{8}$ έχουν αδέλφια.\par
\tithelo Πόσοι έχουν αδέλφια και πόσοι όχι.\par
\possindeo Τα $\frac{5}{8}$ του 96 είναι $\frac{5}{8}\cdot 96=\frac{5\cdot 96}{8}=\frac{480}{8}=60$. Οι υπόλοιποι είναι $96-60=36$.\par
\apantisi 60 έχουν αδέλφια και 36 δεν έχουν.\par
\elegxos Τα υπόλοιπα $\frac{3}{8}$ του 96 είναι $\frac{3\cdot 96}{8}=36$. Συμφωνεί.

\paradeigma{Μεικτοί αριθμοί}
\tixero Οι αριθμοί $2\frac{1}{3}$ και $1\frac{1}{2}$.\par
\tithelo Το γινόμενό τους.\par
\possindeo $2\frac{1}{3}=\frac{7}{3}$ και $1\frac{1}{2}=\frac{3}{2}$. Άρα $\frac{7}{3}\cdot\frac{3}{2}=\frac{7\cdot\cancel{3}}{\cancel{3}\cdot 2}=\frac{7}{2}=3\frac{1}{2}$.\par
\apantisi $3\dfrac{1}{2}$.\par
\elegxos Το $2\frac{1}{3}$ είναι λίγο πάνω από 2, και το $1\frac{1}{2}$ είναι το ενάμισι. Ενάμισι φορά το 2 και κάτι είναι λίγο πάνω από 3. Το $3\frac{1}{2}$ βγάζει νόημα.

\paradeigma{Οι αντίστροφοι}
\tixero Οι αριθμοί $\frac{4}{9}$, 6 και $1\frac{1}{2}$.\par
\tithelo Οι αντίστροφοί τους.\par
\possindeo Γυρίζω το κλάσμα ανάποδα. Για το $\frac{4}{9}$: $\frac{9}{4}$. Για το $6=\frac{6}{1}$: $\frac{1}{6}$. Για το $1\frac{1}{2}=\frac{3}{2}$: $\frac{2}{3}$.\par
\apantisi $\dfrac{9}{4}$, $\dfrac{1}{6}$ και $\dfrac{2}{3}$.\par
\elegxos $\frac{4}{9}\cdot\frac{9}{4}=1$, $6\cdot\frac{1}{6}=1$, $\frac{3}{2}\cdot\frac{2}{3}=1$. Σωστά.

\tmima{Ασκήσεις}

\zoni{Το κατάλαβα}{εφαρμογή του κανόνα}

\ex Βρες τα γινόμενα και απλοποίησε: α) $\frac{2}{3}\cdot\frac{5}{7}$ \quad β) $\frac{3}{4}\cdot\frac{8}{9}$ \quad γ) $\frac{5}{6}\cdot\frac{12}{25}$ \quad δ) $\frac{7}{10}\cdot\frac{20}{21}$ \quad ε) $\frac{4}{9}\cdot\frac{3}{8}$ \quad στ) $\frac{9}{14}\cdot\frac{7}{12}$.

\ex Υπολόγισε: α) $4\cdot\frac{2}{5}$ \quad β) $\frac{3}{7}\cdot 14$ \quad γ) $6\cdot\frac{5}{9}$ \quad δ) $12\cdot\frac{3}{8}$ \quad ε) $\frac{5}{6}\cdot 18$ \quad στ) $100\cdot\frac{7}{10}$.

\ex Βρες: α) το $\frac{1}{2}$ του $\frac{1}{2}$ \quad β) το $\frac{1}{3}$ του $\frac{1}{5}$ \quad γ) τα $\frac{2}{3}$ του $\frac{3}{4}$ \quad δ) τα $\frac{3}{4}$ του $\frac{4}{5}$.

\ex Υπολόγισε τα γινόμενα με μεικτούς: α) $2\frac{1}{3}\cdot 3$ \quad β) $1\frac{1}{2}\cdot 2\frac{2}{3}$ \quad γ) $3\frac{1}{4}\cdot 2$ \quad δ) $2\frac{1}{2}\cdot 1\frac{3}{5}$.

\ex Βρες τον αντίστροφο των αριθμών: $\frac{7}{5}$, 3, $\frac{1}{8}$, 1, $\frac{11}{4}$ και $2\frac{1}{2}$.

\ex Συμπλήρωσε ώστε το γινόμενο να είναι 1: α) $\frac{3}{4}\cdot\square$ \quad β) $\frac{5}{9}\cdot\square$ \quad γ) $6\cdot\square$ \quad δ) $\square\cdot\frac{2}{7}$.

\ex Απλοποίησε πριν πολλαπλασιάσεις: α) $\frac{14}{15}\cdot\frac{25}{28}$ \quad β) $\frac{18}{35}\cdot\frac{21}{27}$ \quad γ) $\frac{8}{9}\cdot\frac{15}{16}$ \quad δ) $\frac{4}{21}\cdot\frac{35}{12}$.

\ex Βρες την τιμή, με προτεραιότητα των πράξεων: α) $\frac{2}{3}+\frac{1}{2}\cdot\frac{3}{4}$ \quad β) $\left(\frac{2}{3}+\frac{1}{2}\right)\cdot\frac{3}{4}$ \quad γ) $\frac{1}{2}\cdot\left(\frac{3}{4}-\frac{1}{6}\right)$.

\ex Σωστό ή λάθος; α) $\frac{2}{3}\cdot\frac{3}{2}=1$ \quad β) $\frac{3}{5}\cdot 4=\frac{12}{20}$ \quad γ) $\frac{1}{2}\cdot\frac{1}{3}=\frac{1}{5}$ \quad δ) Ο αντίστροφος του 0 είναι το 0. \quad ε) $4\cdot\frac{1}{4}=1$ \quad στ) $2\frac{1}{2}\cdot 2=4\frac{1}{2}$.

\zoni{Το χρησιμοποιώ}{προβλήματα από τη ζωή}

\ex \textbf{Η συνταγή.} Η συνταγή της γιαγιάς Ρήνης για 6 άτομα θέλει $\frac{3}{4}$ του κιλού αλεύρι. Η Μυρτώ θέλει να μαγειρέψει για 15 άτομα, δηλαδή για $2\frac{1}{2}$ φορές περισσότερα. Πόσο αλεύρι χρειάζεται;

\ex \textbf{Ο χυμός.} Ένα μπουκάλι έχει $1\frac{1}{2}$ λίτρα χυμό. Ο Νίκος πίνει τα $\frac{2}{5}$ του μπουκαλιού. Πόσα λίτρα ήπιε;

\ex \textbf{Ο κήπος.} Τα $\frac{3}{5}$ του κήπου της γιαγιάς είναι λαχανικά και τα $\frac{2}{3}$ από αυτά είναι ντομάτες. Τι μέρος ολόκληρου του κήπου είναι οι ντομάτες; Αν ο κήπος έχει 300 τ.μ., πόσα τ.μ. έχουν ντομάτες;

\ex \textbf{Το εμβαδό.} Ένα ορθογώνιο έχει μήκος $3\frac{1}{2}$ μ. και πλάτος $2\frac{2}{5}$ μ. Βρες το εμβαδό του.

\ex \textbf{Ο περίπατος.} Ο Βασίλης περπατάει $4\frac{1}{2}$ χλμ. την ώρα. Πόσα χιλιόμετρα κάνει σε $\frac{2}{3}$ της ώρας, δηλαδή σε 40 λεπτά;

\ex \textbf{Η έκπτωση.} Ο Νίκος βρίσκει αθλητικά παπούτσια που κοστίζουν 80 € και έχουν έκπτωση τα $\frac{3}{10}$ της τιμής τους. Πόσα ευρώ είναι η έκπτωση; Πόσο πληρώνει; (Δοκίμασε να βρεις και τα $\frac{7}{10}$ της τιμής.)

\zoni{Το σκέφτομαι}{μια ερώτηση που πάει πιο μακριά}

\ex Ο Στέλιος λέει: «Ο πολλαπλασιασμός κάνει τα πράγματα μεγαλύτερα, άρα το $\frac{1}{2}\cdot\frac{1}{3}$ είναι πάνω από το $\frac{1}{2}$». Εξήγησε με το «μισό του ενός τρίτου» ότι βγαίνει μικρότερο.

\ex Πολλαπλασίασε το 12 επί $\frac{3}{4}$ και επί $\frac{5}{4}$. Πότε ο πολλαπλασιασμός με κλάσμα μικραίνει τον αριθμό και πότε τον μεγαλώνει;

\ex Ποιο κλάσμα πρέπει να πολλαπλασιάσεις με το $\frac{3}{8}$ για να βρεις $\frac{1}{2}$; (Σκέψου τον αντίστροφο του $\frac{3}{8}$.)

\ex Υπολόγισε $\frac{1}{2}\cdot\frac{2}{3}\cdot\frac{3}{4}\cdot\frac{4}{5}\cdot\frac{5}{6}$. Τι παρατηρείς; Μάντεψε πόσο θα βγει το $\frac{1}{2}\cdot\frac{2}{3}\cdot\dots\cdot\frac{9}{10}$.

\ex Γιατί το 0 δεν έχει αντίστροφο; Εξήγησε με την εξίσωση $0\cdot x=1$.

\zoni{Για το σπίτι}{με χαρτί και με το μυαλό}

\ex Πάρε ένα φύλλο χαρτί, δίπλωσέ το σε 3 ίσες σειρές και σε 4 ίσες στήλες. Σκίασε τις 3 από τις 4 στήλες και μετά σκούρυνε 2 από τις 3 σειρές πάνω στη σκιασμένη περιοχή. Δείξε με το φύλλο ότι $\frac{2}{3}\cdot\frac{3}{4}=\frac{1}{2}$.

\ex Βρες το $\frac{1}{2}$ του $\frac{1}{2}$ του $\frac{1}{2}$ του $\frac{1}{2}$ του 80.
